\section{Conclusions}\label{sec:conclusion3}\label{sec:extension_limits}
We have extended our SECRYPT 2026 study of Zone~\cite{zone_secrypt_2026}
from its original native transformations to a broader evaluation of
binary rewriting. The extension makes code--data identification and
supported reference handling explicit, separates the LLM roles of
classification and instruction transformation, and integrates the
existing assembly-wrapping technique on x64 and x86. These additions
address the preparation and evaluation of transformed programs alongside
the transformation itself.

The recovery results demonstrate a substantial improvement on inputs
that were difficult for the earlier pipeline. Working rate rises from
1.2\% to 80.2\% on the selected x64 cohort and from 14.9\% to 25.9\% on
x86. Replacing the retested inputs' rates while retaining historical
weights gives mean configuration rates of 82.6\% and 50.3\%, compared
with 55.6\% and 45.0\% historically. This combined summary places the
recovery results in the wider corpus; it does not represent a fresh
execution of every historical test. The gain is larger on x64, while
the residual x86 cases show that the improvement is architecture- and
input-dependent.

The controlled suites and original payloads provide complementary
evidence. Six annotated programs give known answers for data
identification, and 13 functional programs isolate individual programming
patterns. Their small, simple layouts explain why their high success
rates do not automatically extend to the larger corpus. On the 19
original data-handling inputs, symbolic and LLM proposals yield identical
final binaries in 16 cases and the same 17 successful functional outcomes.
This agreement shows that different intermediate proposals can lead to
the same accepted output under shared handling. It also keeps the
classifier's contribution distinct from that of the downstream pipeline.

The full-corpus campaigns expose the difference between producing a
binary and obtaining a working variant. Wrapping constructs 243 binaries
from 306 planned cases, of which 170 pass their functional tests. Every
constructed x64 wrapping output passes, whereas x86 functionality remains
partial. LLM transformation constructs 350 binaries across 816 conditions,
with 275 successful functional outcomes. These results identify
construction failures as a substantial part of the measured outcome;
reporting only accepted outputs would conceal that part of the process.
The recorded native x86 IBP result remains tied to its tested engine
version, separately from these later campaigns.

Zero-shot and two-shot prompting produce similar aggregate outcomes in
this experiment. The evidence supports no consistent advantage for the
additional examples under the tested settings. Static scans report no
malicious responses for the 207 distinct wrapping outputs and a mean
engine detection rate of 5.06\% across 322 nonempty LLM-output reports.
These observations describe the submitted raw files and complement the
functional results. Overall, the extension demonstrates improved
functional reliability and provides a more complete account of the
added transformation paths by reporting construction, execution,
diversity, and detection together.
